\section*{Cinco principios holográficos de construcción y conservación}
\addcontentsline{toc}{section}{Cinco principios holográficos de construcción y conservación}

Los \emph{Elementos} ofrecen un antecedente preciso para presentar una
arquitectura matemática mediante operaciones declaradas y consecuencias
demostradas. El libro I distingue definiciones, postulados y nociones comunes.
Sus tres primeros postulados permiten trazar un segmento entre puntos,
prolongar un segmento en línea recta y describir un círculo con centro y radio
dados; el cuarto establece la igualdad de los ángulos rectos. El quinto
formula una condición de encuentro: dos rectas cortadas por una transversal
se encuentran, al prolongarse, del lado donde los ángulos interiores suman
menos de dos rectos. Son enunciados sobre los objetos y las construcciones de
esa geometría.\footnote{Euclides, \emph{Elementos}, libro I, postulados 1--5.
Texto griego de Heiberg y traducción de Richard Fitzpatrick, p.~7:
\url{https://farside.ph.utexas.edu/books/Euclid/Elements.pdf}. También pueden
consultarse por separado en la edición de David E. Joyce:
\url{https://mathcs.clarku.edu/~djoyce/elements/bookI/bookI.html}.}

La extrema y media razón aparece en el libro VI como una relación entre un
todo y sus partes: el todo es al segmento mayor como éste al menor. En una
traducción algebraica moderna, para longitudes positivas \(u>v\), se escribe
\((u+v)/u=u/v\); por tanto, la razón \(r=u/v\) satisface
\(r^2=r+1\). Esta escritura explica la proporción transmitida sin convertir
el lenguaje algebraico moderno en el punto de partida de Euclides. En HMT,
la coordenada de autoescala se obtiene desde una incidencia producida por el
calendario operativo; su reconocimiento mediante la misma ecuación pertenece
a una etapa posterior.\footnote{Euclides, \emph{Elementos}, libro VI,
definición 3, en la edición de Joyce:
\url{https://mathcs.clarku.edu/~djoyce/elements/bookVI/defVI3.html}.
La edición de Fitzpatrick imprime este mismo enunciado como definición 2
del libro VI, p.~156. Se conservan aquí ambas numeraciones editoriales.}

La comparación que sigue se sitúa en la organización de la construcción.
HMT declara un soporte aritmético, dos hojas relacionadas, una orientación
trítica y operaciones que prolongan estados conservando sus registros. Un
cilindro representa la información fijada por un prefijo; una recta
euclidiana es otro objeto, con otra definición y otras operaciones. Los
cinco principios de esta sección expresan relaciones efectivas entre los
objetos holográficos. Su número ofrece una exposición reconocible; no
establece una correspondencia término a término con los cinco postulados,
ni una refutación del quinto, ni una sustitución de los axiomas geométricos.

Estos principios tienen estatuto \emph{organizador}: reúnen condiciones y
teoremas que el cuerpo del artículo construye y prueba. Cada uno presenta primero
su sentido, después sus dominios y una identidad que permite comprobarlo.
La secuencia común es
\[
 \mathrm{APP}\longrightarrow\mathrm{TRIT}\longrightarrow\mathrm{TPK}
 \longrightarrow\text{estado enriquecido}
 \longrightarrow\text{estructura discreta conjunta del continuo}.
\]
Las realizaciones numéricas, lineales y geométricas leen después esa
estructura. Esta posición de las realizaciones mantiene visibles el origen
de los coeficientes y la información que cada lectura conserva.

\subsection*{Principio I. Unidad relacional de la generación}
\hypertarget{prinh:unidad}{}

\emph{Cada lectura se construye mediante operaciones tipadas sobre un
mismo soporte, conservando las relaciones que permiten continuarla.} Una
casilla APP admite simultáneamente una lectura aditiva y otra multiplicativa.
El paso a una cifra visible registra una parte de ambas operaciones; el
cociente, la orientación y la hoja conservan las condiciones de su
prolongación. La unidad consiste en esta procedencia relacional, que permite
seguir una misma ejecución a través de sus distintas realizaciones.

El dominio inicial es \(D_9^2\), con \(D_9=\{1,\ldots,9\}\). Sobre él
actúan \(S(i,j)=i+j\) y \(P(i,j)=ij\). Para un entero \(n\), la coordenatización
positiva de residuo y cociente viene dada por
\[
 \rho_9(n)=1+((n-1)\bmod9),\qquad
 q_9(n)=\frac{n-\rho_9(n)}9,\qquad
 n=\rho_9(n)+9q_9(n).
\]
La evaluación levantada de APP conserva
\(\bigl((\rho_9(S),q_9(S)),(\rho_9(P),q_9(P))\bigr)\).
La identidad reconstruye las dos lecturas enteras; por ejemplo, distingue
los valores 10 y 19, que tienen el mismo residuo positivo y cocientes
diferentes. El resultado y su prueba se desarrollan en
\eqref{eq:residuo-cociente} y \eqref{eq:app-levantada}. La representación
por residuos \(0,\ldots,8\), utilizada en otras coordenatizaciones del artículo, conserva
su cociente correspondiente: cambiar de representantes exige cambiar ambos
componentes de la división.

TRIT añade la representación orientada de \(\mathbb F_3\) por
\(\{-1,0,+1\}\), con su división exacta \(n=r+3c\). El selector operativo
\(M_{\rm ph}\) toma los valores \(+1,-1,0\) en las clases de fase
\(\{1,4,7\}\), \(\{2,5,8\}\), \(\{3,6,9\}\), respectivamente.
El transporte cardinal tiene otro dominio y otra función: mueve los
cursores sobre \(X_9=(\mathbb Z/9\mathbb Z)^2\). En el estado enriquecido
\(\widehat X\), TPK compone selección, transporte y actualización de
memoria,
\[
 U_t=\operatorname{Mem}_t\circ\operatorname{Tra}_t
                    \circ\operatorname{Sel}_t.
\]
La definición \eqref{eq:actualizacion} y el desarrollo de
\cref{tpk-int:estado,tpk-int:seleccion} especifican fase, cursores, hojas,
direcciones, residuos, cocientes y registros. El signo de fase y la dirección
cardinal se componen conservando sus tipos distintos.

La autoescala ofrece un ejemplo de generación que puede seguirse sin partir
de su valor final. El calendario \((+1,-1,0)\) produce, con la regla de
cambio y conservación de hoja, el ciclo modal
\((\Sigma,\Pi,\Pi)\). Sus tres aristas determinan
\[
 F_{\rm av}=\begin{pmatrix}0&1\\1&1\end{pmatrix}.
\]
En su realización lineal posterior, el rayo propio positivo normalizado
como \((1,x)^{\mathsf T}\) satisface \(x=\lambda\) y
\(1+x=\lambda x\). Así se obtiene \(x^2-x-1=0\); la raíz positiva se
reconoce como \(\varphi\). El teorema que sigue a \eqref{eq:fav} prueba
esta identificación después de construir las entradas de la matriz.
La relación con la proporción euclidiana es exacta en esta lectura de
autoescala y conserva la diferencia entre sus dos procedimientos de origen.

\subsection*{Principio II. Prolongación orientada y memoria de composición}
\hypertarget{prinh:prolongacion}{}

\emph{Una prolongación admisible añade un acontecimiento al estado y
transporta la memoria anterior mediante la misma operación.} La historia
permite distinguir recorridos que terminan en una posición o una fase
iguales. Este principio proporciona una manera precisa de hablar de
continuación: se conserva el prefijo efectivamente recorrido y se añade
una emisión permitida por el selector y el régimen presentes.

Sea \(X_k\) el conjunto de historias admisibles de longitud \(k\),
construidas desde las semillas del núcleo. Si \(\varepsilon\) es admisible
tras \(h\), la prolongación \(h\varepsilon\) pertenece a \(X_{k+1}\)
y la truncación satisface \(\rho_k(h\varepsilon)=h\). Se retienen la
etiqueta de emisión y las hojas anteriores, incluso cuando dos emisiones
tengan la misma lectura final. En una coordenatización de memoria sobre un módulo
\(M\), la actualización tiene la forma
\(U_\varepsilon(m)=C_\varepsilon m+\kappa_\varepsilon\), con
\(C_\varepsilon\in\operatorname{End}(M)\) y
\(\kappa_\varepsilon\in M\). Si dos tramos son componibles en esa
coordenatización, entonces
\[
 C_{\beta\alpha}=C_\beta C_\alpha,\qquad
 \kappa_{\beta\alpha}=C_\beta\kappa_\alpha+\kappa_\beta,
 \qquad U_{\beta\alpha}=U_\beta U_\alpha.
\]
La prueba de \cref{iv:cont:concatenacion} consiste en sustituir la primera
actualización en la segunda. La memoria se suma después del transporte que
le corresponde; la asociatividad hace independiente el resultado de la
subdivisión del trayecto.

La diferencia entre fase y memoria aparece ya en una lectura elemental.
Para \((w,r)\in\mathbb Z\times C_9\), con
\(\bar r\in\{0,\ldots,8\}\), se tiene
\[
 \sigma_9(w,r)=
 \left(w+\left\lfloor\frac{\bar r+1}{9}\right\rfloor,
       r+1\pmod9\right),\qquad
 \sigma_9^9(w,r)=(w+1,r).
\]
La segunda igualdad cuenta exactamente un cruce de 8 a 0. Su primera
coordenada registra el avance que la segunda ya no distingue. En el
calendario de doce ventanas, la sucesión exacta
\(0\to C_{12}\to C_{108}\to C_9\to0\) conserva el acarreo entre
ventanas. No se escinde: \(C_{108}\) tiene exponente 108 y
\(C_{12}\times C_9\), exponente 36. La prueba y la ley del acarreo están
en \cref{prop:k-extension}.

El retorno de fase adquiere así un significado operativo: permite
comparar lecturas realizadas en una misma posición del calendario mientras
el estado conserva lo sucedido entre ellas. La holonomía \(\Gamma_9\)
se obtiene de la dinámica enriquecida; \(K_{\rm ph}\) es su representación
reducida posterior. El registro dodecafásico entero \(K\), que utiliza
además orientación e incidencias, tiene también su construcción y su dominio
propios. La igualdad de una fase, de un residuo o de una coordenada expresada
no identifica por sí sola las historias. Ésta es la información de partida
que se conserva al pasar al refinamiento.

\subsection*{Principio III. Compatibilidad del refinamiento y realización}
\hypertarget{prinh:refinamiento}{}

\emph{Refinar consiste en añadir resolución compatible con el prefijo,
manteniendo las relaciones de truncación y las masas de sus prolongaciones.}
La lectura a una escala más fina precisa qué aconteció dentro de un
cilindro; al olvidar ese detalle se recupera la descripción anterior.
La realización de un valor surge cuando los cilindros de su lectura se
contraen de manera compatible. La historia que determina esos cilindros
permanece en el dominio de la construcción.

En el árbol de \cref{iv:cont:medida}, el conjunto inicial es finito y no
vacío, y cada historia tiene un número finito y positivo de sucesores
admisibles. Una masa inicial positiva de suma uno y probabilidades locales
\(p(\varepsilon\mid h)>0\), de suma uno en cada predecesor, producen
\[
 \mu(h\varepsilon)=\mu(h)p(\varepsilon\mid h),\qquad
 \sum_{\rho_k(y)=h}\mu(y)=\mu(h).
\]
Por tanto, en la realización posterior
\(H_k=L^2(X_k,\mu_k)\), el levantamiento
\((J_kf)(y)=f(\rho_k(y))\) es isométrico. Su adjunto es la esperanza
condicional
\[
 (E_kg)(h)=\sum_{\rho_k(y)=h}
                  \frac{\mu(y)}{\mu(h)}g(y),
 \qquad E_kJ_k=I.
\]
Agrupar la norma cuadrada por predecesores demuestra la isometría y agrupar el
producto interno demuestra la fórmula del adjunto
(\cref{iv:cont:esperanza}). El levantamiento de una lectura y su retorno
al nivel anterior forman así un par de operaciones explícitas.

Las hipótesis de ramificación finita y prolongación no vacía proporcionan
historias infinitas compatibles. Sus cilindros reciben las masas anteriores,
que determinan una medida sobre la sigma-álgebra cilíndrica. La elección de
\(p\) pertenece al lector declarado: el caso de sucesores equiprobables da
\(p=1/d(h)\), sin imponer esa medida a cualquier otra lectura. Esta
construcción conserva multiplicidades de historias, incluidas las que
comparten un valor realizado. La prolongabilidad se refiere al árbol
admisible construido; una restricción adicional de rutas necesita conservar
expresamente esa propiedad.

Para la lectura numérica regional, cada bloque
\(b\in\mathbb F_3^6\) tiene un índice entero
\(\nu(b)=\sum_{j=1}^6 b_j3^{6-j}\), comprendido entre 0 y 728.
El lector fija \(N_0=m\) y produce
\[
 N_{n+1}=729N_n+\nu(b_{n+1}),\qquad
 I_n=\left[\frac{N_n}{729^n},\frac{N_n+1}{729^n}\right].
\]
Se cumple \(I_{n+1}\subseteq I_n\) y
\(\operatorname{diam}(I_n)=729^{-n}\). Los extremos izquierdos forman
una sucesión racional de Cauchy; su clase define primero la realización
interna de la lectura. Identificada posteriormente con la completación
usual de los racionales, corresponde al único punto de los intervalos
cerrados anidados. El teorema posterior a \eqref{eq:cilindro} demuestra
esta afirmación y trata por separado las convenciones de frontera. Para
los caracteres regionales, el lector efectivo de cada prefijo decimal
utiliza las horquillas calculables y la separación demostrada respecto de
sus fronteras, como precisa \cref{gen:prefijos-regionales}. Esa profundidad
arbitraria conserva el cuantificador «para cada longitud finita»; la
existencia de un límite por compatibilidad y contracción es una afirmación
distinta y no aporta, por sí sola, un algoritmo para cualquier historia
abstracta. Las historias que representan el mismo punto conservan sus
registros distintos.

\subsection*{Principio IV. Doble lectura correlativa y transporte de marco}
\hypertarget{prinh:doble}{}

\emph{Las lecturas numérica e incidencial proceden del mismo estado
enriquecido y conservan sus compatibilidades al prolongar y cambiar de
marco.} La primera precisa una realización mediante intervalos cada vez
más estrechos. La segunda retiene relaciones entre posiciones, signos y
fronteras. El vínculo holográfico consiste en sus mapas de origen y
transporte comunes, que permiten comparar qué información expresa cada una.

Sea \(X_n^{\rm cal}\) el dominio de prefijos admisibles que incluyen
su marco inicial y los transportes sucesivos de marco. La lectura aritmética
utiliza los bloques de ese prefijo; la incidencial utiliza las palabras
\(w_j\in\mathbb F_3^6\) y los marcos \(f_j\). Para una matriz
\(A\in M_6(\mathbb F_3)\), el lector lineal
\[
 \gamma_A:\mathbb F_3^6\longrightarrow\mathbb F_3^{12},
 \qquad \gamma_A(w)=(w,wA)
\]
actúa sobre palabras escritas como filas. Su especialización en la matriz
Paley--Witt construida en \cref{exc:paley-codigo} produce el código
ternario y sus incidencias. La matriz se obtiene en esa construcción
finita; un número real terminal no selecciona sus entradas. El orden
\(w_6\to w_{12}\to w_{18}\to w_{24}\to w_{30}\to R_{36}\to G_9\)
mantiene las transformaciones y memorias utilizadas antes de cada lectura.

La compatibilidad puede comprobarse sobre el propio operador. Para
\(f\in\operatorname{GL}_6(\mathbb F_3)\), sea
\(A^f=fAf^{-1}\). Entonces
\[
 \gamma_{A^f}(wf^{-1})
      =\gamma_A(w)(f^{-1}\oplus f^{-1}),
\]
pues ambos lados son \((wf^{-1},wAf^{-1})\). En particular, la lectura
del prefijo
\[
 Q_n(x)=\bigl(f_j,
       \gamma_{f_jA_Wf_j^{-1}}(w_j)\bigr)_{0\leq j<n}
\]
satisface \(r_{n,m}Q_n=Q_m p_{n,m}\), donde \(p_{n,m}\) trunca
historias y \(r_{n,m}\) trunca registros. La causalidad de la
actualización de los marcos garantiza que ambos procedimientos retienen
exactamente los mismos términos. El resultado
\cref{exc:limite-inverso} construye su paso al límite. La covarianza lineal
es válida para el cambio de marco escrito; conservar además un soporte
coordenado o una métrica exige transportar esa estructura, o restringir los
cambios de marco al grupo que la preserva.

La construcción correlativa del continuo extiende esta compatibilidad al
transporte por hojas, al balance orientado, a la incidencia ternaria, al
complejo rectangular de puertos y a la memoria celular con sus rellenos y
líneas determinantes. Estos aspectos se construyen sobre las mismas
historias, no sobre cinco dominios elegidos después. Las identidades de
concatenación, naturalidad y borde aseguran su ensamblaje en
\cref{iv:rectangular-natural,iv:cont:memoria-natural,iv:cont:ensamblaje,iv:cont:mapa-correlativo}.
Los cinco principios expositivos de esta apertura no son una nueva
partición de esas cinco construcciones consustanciales. La recuperación
de una palabra mediante la primera proyección de \(\gamma_A\), por su
parte, se refiere a esa palabra: la recuperación de la historia enriquecida
completa exige conservar los demás registros que la distinguen.

\subsection*{Principio V. Conservación y recuperación del registro completo}
\hypertarget{prinh:conservacion}{}

\emph{La información recuperable se conserva conjuntamente en el estado
que continúa, en los registros producidos y en sus correlaciones.} Una
lectura visible puede contraerse mientras el registro completo conserva una
representación isométrica del estado de partida. Este principio permite
formular una reversibilidad verificable: se declara qué se almacena, en qué
espacio reside y mediante qué operador se reconstruye.

El transporte de historias levantadas de \cref{iv:cont:transporte} actúa
entre los espacios \(H_k=\ell^2(\widehat X_k)\) mediante
\[
 \mathcal U_ke_{\widehat h}=
 \sum_{\varepsilon\ {
m admisible}}
 \sqrt{p(\varepsilon\mid\widehat h)}\,
 e_{\widehat h\varepsilon},\qquad
 \mathcal U_k^*\mathcal U_k=I.
\]
La segunda igualdad procede de la normalización local y de que predecesores
distintos tienen sucesores con etiquetas distintas. Si \(P_k\) lee la hoja
\(\Sigma\) y \(Q_k=I-P_k\) la hoja \(\Pi\), las operaciones
\[
 \begin{aligned}
 A_k&=P_{k+1}\mathcal U_kP_k+Q_{k+1}\mathcal U_kQ_k,\\
 \eta_k&=Q_{k+1}\mathcal U_kP_k+P_{k+1}\mathcal U_kQ_k
 \end{aligned}
 \qquad\text{satisfacen}\qquad A_k^*A_k+\eta_k^*\eta_k=I.
\]
La identidad de Gram de \cref{iv:cont:gram} se demuestra al anular los
productos de proyectores complementarios. Distingue permanecer en una
hoja de cambiar de hoja; el valor trítico cero continúa siendo una
orientación y no se convierte en una tercera hoja.

Escribamos \(F_0=I\), \(F_{n+1}=A_nF_n\). La telescopía de la
identidad anterior da, para todo horizonte \(N\),
\[
 \|x\|^2=\|F_Nx\|^2+
                 \sum_{n=0}^{N-1}\|\eta_nF_nx\|^2.
\]
El primer término es el residuo terminal que todavía continúa; los otros
son las salidas registradas. El teorema \cref{iv:cont:terminal} demuestra
que \(F_N^*F_N\) converge fuertemente a un operador positivo
\(T_\infty\). Se obtiene, por tanto, el registro completo
\[
 \mathcal V:H_0\longrightarrow H_0\oplus
                     \bigoplus_{n\geq0}H_{n+1},\qquad
 \mathcal Vx=\bigl(T_\infty^{1/2}x,
                         (\eta_nF_nx)_{n\geq0}\bigr),
 \qquad\mathcal V^*\mathcal V=I.
\]
En efecto, la suma de las normas cuadradas es \(\|x\|^2\); la
polarización da la última identidad. Su imagen \(\mathcal M\) es
cerrada y la inversa sobre ella es \(\mathcal V^*\). El terminal se
conserva hasta demostrar su extinción; \(T_\infty\) no necesita ser
un proyector ni se supone la convergencia de \(F_Nx\).

La conservación alcanza también los operadores que leen esa información.
Para un operador acotado \(B\) sobre \(H_0\), el transporte
\(B^{\mathcal V}=\mathcal VB\mathcal V^*|_{\mathcal M}\) conserva
productos, adjuntos y productos internos:
\[
 (B_1B_2)^{\mathcal V}=B_1^{\mathcal V}B_2^{\mathcal V},\qquad
 \langle\mathcal Vx,B^{\mathcal V}\mathcal Vy\rangle
                       =\langle x,By\rangle.
\]
Insertar \(\mathcal V^*\mathcal V=I\) prueba ambas fórmulas. En el
lector nonádico de \cref{lib01:isometria}, las mismas identidades se
realizan con \(T_n=(8I+C_n)/9\) y
\(D_n=\sqrt8(C_n-I)/9\), para unitarias \(C_n\) del portador
declarado. La composición mantiene su estatuto de lector y no identifica
automáticamente cada \(C_n\) con la holonomía generadora \(\Gamma_9\).
La recuperación preserva así las correlaciones del registro y los
observables transportados, dentro de la imagen precisa donde actúa.

El registro \(K\) permite ver esta distinción en un objeto finito. Su
codificación \(\kappa(K)\), en base \(B=1000\) y longitud doce, es
\[
 \kappa(K)=\frac{N(K)}{B^{12}-1},\qquad
 N(K)=\sum_{j=1}^{12}K_jB^{12-j},\qquad
 K\in\{0,\ldots,999\}^{12}.
\]
Multiplicar por \(B^{12}-1\) y efectuar doce divisiones euclídeas
recupera el registro exactamente, incluidos sus bloques iniciales nulos.
La proposición \cref{lib01:inversion-kappa} prueba además que un error
absoluto menor que \(1/[2(B^{12}-1)]\) permite recuperar el mismo entero
mediante redondeo. La reversibilidad se refiere al registro en su coordenatización y
con sus metadatos; no identifica a éste con toda la historia TPK. La
\emph{constante holográfica de acoplamiento aritmético-geométrico} designa
en este artículo el objeto \(K\) tipado con sus lectores y relaciones. Su
representación procede del estado enriquecido y participa después en la
clausura de alfa: no se utiliza como entrada generadora de \(\pi\), ni
su denominación convierte cada lectura escalar en un invariante universal.

\subsection*{Una lectura conjunta de los cinco principios}

La secuencia ofrece un criterio positivo para recorrer el artículo. La unidad
relacional fija qué se opera; la prolongación conserva lo recorrido; el
refinamiento organiza la resolución; la doble lectura mantiene juntos valor
e incidencia; la recuperación identifica las condiciones exactas para
reconstruir el registro y sus observables. Las demostraciones se enlazan
porque actúan sobre objetos con procedencia común y respetan sus mapas de
transporte. Allí donde una lectura olvida información, su dominio y su
fibra de pérdida permanecen definidos.

La extrema y media razón histórica muestra cómo una relación entre el todo
y sus partes puede determinar una proporción. La construcción holográfica
desarrollada aquí estudia, además, relaciones conservadas cuando cambian la
resolución, el marco y la distribución de información entre continuidad y
memoria. La ecuación de autoescala establece un enlace matemático preciso
con aquella proporción; las identidades de concatenación, refinamiento e
isometría proporcionan el contenido adicional que este artículo desarrolla.
La comparación histórica conserva así su función orientadora, mientras
cada conclusión holográfica descansa en sus operaciones y pruebas propias.

\subsection*{Fuentes primarias de la comparación histórica}

\begin{enumerate}[label=\arabic*.]
 \item Euclides, \emph{Elementos}, libro I, postulados 1--5. Texto
 griego establecido por J. L. Heiberg, con traducción de Richard Fitzpatrick,
 edición corregida de 2008, p.~7; libro VI, definición 2 en esta edición,
 p.~156. \url{https://farside.ph.utexas.edu/books/Euclid/Elements.pdf}.
 \item Euclides, \emph{Elementos}, edición digital de David E. Joyce,
 Clark University: libro I, postulados 1--5,
 \url{https://mathcs.clarku.edu/~djoyce/elements/bookI/bookI.html};
 libro VI, definición 3,
 \url{https://mathcs.clarku.edu/~djoyce/elements/bookVI/defVI3.html}.
 La diferencia de numeración del libro VI se conserva para facilitar el
 cotejo entre ediciones.
\end{enumerate}
